\scanpage{238}
\mspage{112}
% archive sheet number: 237

$\delta_0^X,\delta_1^X$ (qui en sont des sections) sont dans $W_A$, comme $\delta_0^X$ est mono on a donc $\delta_0^X\in\mathrm{TC}$, donc ($\mathrm{TC}$ étant stable par \unsure{coch.\textsuperscript{t}|coch\textsuperscript{t}|coch\textsuperscript{gt}} de base) $\beta\in W_A$, donc enfin $p_f$ (qui est une rétraction de $Z(f)$ sur $Y$) est aussi dans $W_A$.

Il reste à voir que $i_f$ est mono, ce qui se voit sur le diagramme plus complet (cf $\overline{\underline{\mathrm{XII}}}$ p. 230\deleted{\unsure{)}}) que j'ai sous les yeux : $i_f$ apparaît comme composée de deux mono. $X\longrightarrow Y\amalg X$ et $Y\amalg X\longrightarrow Z(f)$

\inserted{\unsure{Le Th.|le th.} 8.4. est démontré.}

\unsure{.....|....} Je vais quand même isoler l'argument formel \unsure{supplémentaire} (que j'ai utilisé surtout sous forme duale) dans un

\uline{Lemme} 8.7. Soit $\unsure{\mathcal{M}|M}$ une catégorie, $W$ un localiseur modérément saturé dans $\mathcal{M}$, $C$ une partie stable par composition et \deleted{\unsure{contenant les}} \inserted{par cochangt.} de base,

On se propose, pour $X$ fixé dans $\mathcal{M}$, de trouver une factorisation fonctorielle de \unsure{tte|toute} flèche $f:X\unsure{\mathrel{\text{---}}|\longrightarrow}Y$ (i.e. du morphisme structural d'un $Y$ dans $X\backslash\mathcal{M}$) en $p_f i_f$, avec $i_f\in C$, $p_f\in W$. On suppose \deleted{que} pour ceci qu'on a
